
\documentclass[aps,prd,onecolumn,notitlepage,nofootinbib]{revtex4-2}
%%%%%%%%%%%%%%%%%%%%%%%%%%%%%%%%%%%%%%%%%%%%%%%%%%%%%%%%%%%%%%%%%%%%%%%%%%%%%%%%%%%%%%%%%%%%%%%%%%%%%%%%%%%%%%%%%%%%%%%%%%%%%%%%%%%%%%%%%%%%%%%%%%%%%%%%%%%%%%%%%%%%%%%%%%%%%%%%%%%%%%%%%%%%%%%%%%%%%%%%%%%%%%%%%%%%%%%%%%%%%%%%%%%%%%%%%%%%%%%%%%%%%%%%%%%%
\usepackage{amsfonts}
\usepackage{eurosym}
\usepackage{amsmath,amssymb}
\usepackage{physics}
\usepackage{enumitem}

\setcounter{MaxMatrixCols}{10}
%TCIDATA{OutputFilter=Latex.dll}
%TCIDATA{Version=5.50.0.2960}
%TCIDATA{<META NAME="SaveForMode" CONTENT="1">}
%TCIDATA{BibliographyScheme=Manual}
%TCIDATA{LastRevised=Saturday, June 28, 2025 17:22:17}
%TCIDATA{<META NAME="GraphicsSave" CONTENT="32">}

\input{tcilatex}
\begin{document}

\title{Pfaffian Compound Matrices and Wedge Product Transformations}
\author{}
\date{}
\maketitle

\section*{Pfaffian Compound Matrices and Wedge Products}

\subsection*{1. Pfaffian Compound Matrix in a Transformation Law}

Let $V \cong \mathbb{R}^4$ and $\{e_1, e_2, e_3, e_4\}$ be a basis for $V$.
Consider a 2-form: 
\begin{equation}
\omega = \sum_{i<j} A_{ij} e_i \wedge e_j \in \Lambda^2 V
\end{equation}
where $A$ is a $4 \times 4$ skew-symmetric matrix.

Let $T: V \to V$ be a linear transformation represented by $M \in \mathrm{GL}%
(4)$. Then $T$ induces a map $\Lambda^2 T: \Lambda^2 V \to \Lambda^2 V$,
with 
\begin{equation}
\Lambda^2 T(e_i \wedge e_j) = (T e_i) \wedge (T e_j)
\end{equation}

The matrix $A$ transforms as: 
\begin{equation}
A\mapsto M^{\prime T}AM
\end{equation}

We can write $\omega^k \in \Lambda^{2k} V$ as: 
\begin{equation}
\omega^k = \sum_I \alpha_I \, e_{i_1} \wedge \cdots \wedge e_{i_{2k}}, \quad 
\text{where } \alpha_I = \func{Pf}(A_I)
\end{equation}

Under the basis change $M$, the coefficients transform as: 
\begin{equation}
\vec{\alpha}^{\prime }= \mathrm{Pf}_k(M)^T \cdot \vec{\alpha}
\end{equation}

This parallels the role of compound matrices in determinant expansions.

\subsection*{2. Example: $4 \times 4$ Case for $\mathrm{Pf}_1(A)$}

Let: 
\begin{equation}
A = 
\begin{pmatrix}
0 & a & b & c \\ 
-a & 0 & d & e \\ 
-b & -d & 0 & f \\ 
-c & -e & -f & 0%
\end{pmatrix}%
\end{equation}

The 6 basis elements of $\Lambda^2 V$ are: 
\begin{align*}
e_{12}, e_{13}, e_{14}, e_{23}, e_{24}, e_{34}
\end{align*}

Label the Pfaffian compound matrix $\mathrm{Pf}_1(A)$ with rows/columns
indexed by these pairs. The diagonal entries are: 
\begin{align*}
\func{Pf}(A_{\{1,2\}}) &= a \\
\func{Pf}(A_{\{1,3\}}) &= b \\
\func{Pf}(A_{\{1,4\}}) &= c \\
\func{Pf}(A_{\{2,3\}}) &= d \\
\func{Pf}(A_{\{2,4\}}) &= e \\
\func{Pf}(A_{\{3,4\}}) &= f
\end{align*}

An off-diagonal entry: 
\begin{align*}
[\mathrm{Pf}_1(A)]_{\{1,2\},\{3,4\}} &= \func{Pf}(A_{\{1,2,3,4\}}) = af - be
+ cd
\end{align*}

Hence, the partial structure of $\mathrm{Pf}_1(A)$ is: 
\begin{equation}
\begin{pmatrix}
a & * & * & * & * & af - be + cd \\* 
& b & * & * & * & * \\* 
& * & c & * & * & * \\* 
& * & * & d & * & * \\* 
& * & * & * & e & * \\ 
af - be + cd & * & * & * & * & f%
\end{pmatrix}%
\end{equation}

The remaining off-diagonal entries are Pfaffians of $4 \times 4$ submatrices
formed by union of index pairs.

\section*{Summary}

\begin{itemize}
\item Under a basis change $M$, the coefficients of $\omega ^{k}$ in $%
\Lambda ^{2k}V$ transform via the Pfaffian compound matrix $\mathrm{Pf}%
_{k}(M)$.

\item For $V=\mathbb{R}^{4}$, $\mathrm{Pf}_{1}(A)\in \mathbb{R}^{6\times 6}$
has diagonal entries given by $2\times 2$ Pfaffians, and off-diagonal
entries built from $4\times 4$ Pfaffians.

\item This structure governs the transformation of even-grade wedge powers
and relates to Pl\"{u}cker coordinates.
\end{itemize}

\section{Goal}

Let $A$ be a skew-symmetric $2n\times 2n$ matrix, and define the 2-form: 
\begin{equation*}
\omega =\sum_{i<j}A_{ij}\,e_{i}\wedge e_{j}\in \Lambda ^{2}V
\end{equation*}%
We want to express $\omega ^{2}\in \Lambda ^{4}V$ in terms of \textbf{%
Pfaffians of $4\times 4$ principal submatrices} of $A$\ the so-called \emph{%
compound Pfaffian matrix} entries.

\section{Step-by-Step Breakdown}

We work in a $2n$-dimensional space $V$ with basis $\{e_1, \dots, e_{2n}\}$.
Fix $n \geq 2$, so $\omega^2 \in \Lambda^4 V$.

\subsection{Quadratic Expansion of $\protect\omega$}

\begin{equation*}
\omega = \sum_{i<j} A_{ij} \, e_i \wedge e_j
\end{equation*}
Then: 
\begin{equation*}
\omega^2 = \sum_{i<j} \sum_{k<\ell} A_{ij} A_{k\ell} (e_i \wedge e_j \wedge
e_k \wedge e_\ell)
\end{equation*}
Only terms with disjoint index pairs $(i<j), (k<\ell)$ contribute. The sum
is over unordered \emph{partitions} of 4 distinct indices into two unordered
pairs.

\subsection{Basis for $\Lambda^4 V$}

All wedge monomials $e_i \wedge e_j \wedge e_k \wedge e_\ell$ with $i < j <
k < \ell$ form a basis of $\Lambda^4 V$. Each such 4-subset $I =
\{i,j,k,\ell\}$ defines a $4 \times 4$ principal submatrix $A_I$.

\begin{equation*}
\omega^2 = \sum_{\substack{ I \subset \{1, \dots, 2n\}  \\ |I| = 4}} \text{Pf%
}(A_I) \cdot \beta_I
\end{equation*}
where $\beta_I = e_{i_1} \wedge e_{i_2} \wedge e_{i_3} \wedge e_{i_4}$ and $%
\text{Pf}(A_I)$ is the Pfaffian of the $4 \times 4$ submatrix $A_I$.

\section{Example: $n = 3$, $\dim V = 6$}

Let $I = \{1,2,3,4\}$. The submatrix $A_I$ is: 
\begin{equation*}
A_I = 
\begin{bmatrix}
0 & A_{12} & A_{13} & A_{14} \\ 
-A_{12} & 0 & A_{23} & A_{24} \\ 
-A_{13} & -A_{23} & 0 & A_{34} \\ 
-A_{14} & -A_{24} & -A_{34} & 0%
\end{bmatrix}%
\end{equation*}
The Pfaffian of $A_I$ is: 
\begin{equation*}
\text{Pf}(A_I) = A_{12} A_{34} - A_{13} A_{24} + A_{14} A_{23}
\end{equation*}
So: 
\begin{equation*}
\text{Contribution to } \omega^2 \text{ from } I = \text{Pf}(A_I) \cdot e_1
\wedge e_2 \wedge e_3 \wedge e_4
\end{equation*}

\section{Connection to Compound Pfaffian Matrix}

This yields: 
\begin{equation*}
\omega ^{2}=\sum_{I\in \binom{2n}{4}}\text{Pf}(A_{I})\cdot e_{I}
\end{equation*}%
where $e_{I}=e_{i_{1}}\wedge e_{i_{2}}\wedge e_{i_{3}}\wedge e_{i_{4}}$ and $%
\text{Pf}(A_{I})$ is the entry of the second compound Pfaffian matrix $\text{%
Pf}_{2}(A)$.

\section{Final Summary}

\begin{equation*}
\backslash omega\symbol{94}2=\backslash sum\_\{I\backslash subset\backslash
\{1,\backslash dots,2n\backslash \},\backslash ,|I|=4\}\backslash
text\{Pf\}(A\_I)\backslash cdote\_I
\end{equation*}%
This is the coordinate expansion of $\omega ^{2}$ in $\Lambda ^{4}V$, with
coefficients given by $4\times 4$ Pfaffian minors of $A$\ the entries of the
compound Pfaffian matrix.

\section*{1. Introduction}

Compound Pfaffians arise naturally in the study of induced transformations
on the even subalgebra of the exterior algebra, particularly in the context
of linear transformations on the 2-forms (grade 2 elements). The
transformation of elements such as $\omega^{\wedge k} \in \Lambda^{2k} V $
under linear maps can be expressed in terms of compound Pfaffians, which are
analogues of compound determinants for skew-symmetric matrices.

\section*{2. Setup: Linear Transformations and Exterior Algebra}

Let $V$ be an $n$-dimensional vector space over a field $\mathbb{F}$, and
let $\Lambda ^{2}V\subset \Lambda ^{\text{even}}(V)$ denote the grade-2
subspace.

Given a linear transformation $T: V \to V $, it induces transformations on $%
\Lambda^2 V $ (denoted $\Lambda^2 T $) and on all higher even subspaces $%
\Lambda^{2k} V $.

If $\omega \in \Lambda^2 V $ corresponds to a skew-symmetric matrix $A $,
then powers like $\omega^{\wedge k} \in \Lambda^{2k} V $ expand using
compound Pfaffians.

\section*{3. Example in Dimension 4}

Let $V = \mathbb{R}^4 $, with basis $\{e_1, e_2, e_3, e_4\} $, and define a
2-form: 
\begin{equation}
\omega = \sum_{1 \le i < j \le 4} A_{ij} \, e_i \wedge e_j,
\end{equation}
with $A $ a skew-symmetric $4 \times 4 $ matrix.

Then the wedge square is: 
\begin{equation}
\omega \wedge \omega = 2 \cdot \func{Pf}(A) \cdot e_1 \wedge e_2 \wedge e_3
\wedge e_4.
\end{equation}

\section*{4. Transformation Properties}

Under a linear transformation $T $ with matrix $M $, we have: 
\begin{equation}
A \mapsto M A M^T,
\end{equation}
and the Pfaffian transforms as: 
\begin{equation}
\func{Pf}(M A M^T) = \det(M) \cdot \func{Pf}(A).
\end{equation}

\section*{5. Higher Dimension Example: $\mathbb{R}^6 $}

Let $V = \mathbb{R}^6 $. A 2-form is: 
\begin{equation}
\omega = \sum_{1 \le i < j \le 6} A_{ij} \, e_i \wedge e_j,
\end{equation}
and $\omega^{\wedge 3} \in \Lambda^6 \mathbb{R}^6 $ is: 
\begin{equation}
\omega^{\wedge 3} = 6! \cdot \func{Pf}(A) \cdot e_1 \wedge e_2 \wedge e_3
\wedge e_4 \wedge e_5 \wedge e_6.
\end{equation}

The Pfaffian of a $6 \times 6 $ skew-symmetric matrix is: 
\begin{equation}
\func{Pf}(A) = \sum_{\pi} \func{sgn}(\pi) A_{\pi(1)\pi(2)} A_{\pi(3)\pi(4)}
A_{\pi(5)\pi(6)},
\end{equation}
where the sum is over partitions into disjoint unordered pairs.

\section*{6. Specific Term Computation}

Assume only: 
\begin{equation}
A_{12} = a, \quad A_{34} = b, \quad A_{56} = c,
\end{equation}
are nonzero. Then: 
\begin{equation}
\omega = a \, e_1 \wedge e_2 + b \, e_3 \wedge e_4 + c \, e_5 \wedge e_6,
\end{equation}
and: 
\begin{equation}
\omega^{\wedge 3} = 6abc \cdot e_1 \wedge e_2 \wedge e_3 \wedge e_4 \wedge
e_5 \wedge e_6.
\end{equation}

So: 
\begin{equation}
\func{Pf}(A) = abc.
\end{equation}

This coefficient is a compound Pfaffian: a product of Pfaffians from three $%
2 \times 2 $ blocks.

\section*{7. General Interpretation}

Each term in $\omega ^{\wedge k}\in \Lambda ^{2k}V$ corresponds to a
compound Pfaffian $\func{Pf}(A_{I})$, where $A_{I}$ is a $2k\times 2k$
principal submatrix. Under transformations $A\mapsto MAM^{T}$, these
compound Pfaffians transform as weighted sums of Pfaffians on new index
sets, mirroring the transformation properties of compound determinants.

\section*{Pfaffian Compound Matrices and Wedge Products}

\subsection*{1. Pfaffian Compound Matrix in a Transformation Law}

Let $V \cong \mathbb{R}^4$ and $\{e_1, e_2, e_3, e_4\}$ be a basis for $V$.
Consider a 2-form: 
\begin{equation}
\omega = \sum_{i<j} A_{ij} e_i \wedge e_j \in \Lambda^2 V
\end{equation}
where $A$ is a $4 \times 4$ skew-symmetric matrix.

Let $T:V\rightarrow V$ be a linear transformation represented by $M\in 
\mathrm{GL}(4)$. Then $M$ induces a map $\Lambda ^{2}M:\Lambda
^{2}V\rightarrow \Lambda ^{2}V$, with 
\begin{equation}
\Lambda ^{2}M(e_{i}\wedge e_{j})=(Me_{i})\wedge (Me_{j})
\end{equation}

The matrix $A$ transforms as: 
\begin{equation}
A\prime \mapsto M^{T}AM
\end{equation}

We can write $\omega ^{k}\in \Lambda ^{2k}V$ as: 
\begin{equation}
\omega ^{k}=\sum_{I}\alpha _{I}\,e_{i_{1}}\wedge \cdots \wedge
e_{i_{2k}},\quad \text{where }\alpha _{I}=\func{Pf}(A\prime _{I})
\end{equation}

Under the basis change $M$, the coefficients transform as: 
\begin{equation}
\vec{\alpha}\prime ^{\prime }=\mathrm{Pf}_{k}(M^{T}AM)^{T}\cdot \vec{\alpha}
\end{equation}

This parallels the role of compound matrices in determinant expansions.

\subsection*{2. Example: $4 \times 4$ Case for $\mathrm{Pf}_1(A)$}

Let: 
\begin{equation}
A = 
\begin{pmatrix}
0 & a & b & c \\ 
-a & 0 & d & e \\ 
-b & -d & 0 & f \\ 
-c & -e & -f & 0%
\end{pmatrix}%
\end{equation}

The 6 basis elements of $\Lambda^2 V$ are: 
\begin{align*}
e_{12}, e_{13}, e_{14}, e_{23}, e_{24}, e_{34}
\end{align*}

Label the Pfaffian compound matrix $\mathrm{Pf}_1(A)$ with rows/columns
indexed by these pairs. The diagonal entries are: 
\begin{align*}
\func{Pf}(A_{\{1,2\}}) &= a \\
\func{Pf}(A_{\{1,3\}}) &= b \\
\func{Pf}(A_{\{1,4\}}) &= c \\
\func{Pf}(A_{\{2,3\}}) &= d \\
\func{Pf}(A_{\{2,4\}}) &= e \\
\func{Pf}(A_{\{3,4\}}) &= f
\end{align*}

An off-diagonal entry: 
\begin{align*}
[\mathrm{Pf}_1(A)]_{\{1,2\},\{3,4\}} &= \func{Pf}(A_{\{1,2,3,4\}}) = af - be
+ cd
\end{align*}

Hence, the partial structure of $\mathrm{Pf}_{1}(A)$ is: 
\begin{equation}
\begin{pmatrix}
a & \ast & \ast & \ast & \ast & af-be+cd \\* 
\ast & b & \ast & \ast & \ast & \ast \\* 
\ast & \ast & c & \ast & \ast & \ast \\* 
\ast & \ast & \ast & d & \ast & \ast \\* 
\ast & \ast & \ast & \ast & e & \ast \\ 
af-be+cd & \ast & \ast & \ast & \ast & f%
\end{pmatrix}%
\end{equation}

The remaining off-diagonal entries are Pfaffians of $4 \times 4$ submatrices
formed by union of index pairs.

\section*{Summary}

\begin{itemize}
\item Under a basis change $M$, the coefficients of $\omega ^{k}$ in $%
\Lambda ^{2k}V$ transform via the Pfaffian compound matrix $\mathrm{Pf}%
_{k}(M)$.

\item For $V=\mathbb{R}^{4}$, $\mathrm{Pf}_{1}(A)\in \mathbb{R}^{6\times 6}$
has diagonal entries given by $2\times 2$ Pfaffians, and off-diagonal
entries built from $4\times 4$ Pfaffians.

\item This structure governs the transformation of even-grade wedge powers
and relates to Pl\"{u}cker coordinates.
\end{itemize}

\section*{Pfaffian Compound Matrices and Wedge Products}

\subsection*{1. Pfaffian Compound Matrix in a Transformation Law}

Let $V \cong \mathbb{R}^4$ and $\{e_1, e_2, e_3, e_4\}$ be a basis for $V$.
Consider a 2-form: 
\begin{equation}
\omega = \sum_{i<j} A_{ij} e_i \wedge e_j \in \Lambda^2 V
\end{equation}
where $A$ is a $4 \times 4$ skew-symmetric matrix.

Let $T:V\rightarrow V$ be a linear transformation represented by $M\in 
\mathrm{GL}(4)$. Then $T$ induces a map $\Lambda ^{2}M:\Lambda
^{2}V\rightarrow \Lambda ^{2}V$, with 
\begin{equation}
\Lambda ^{2}T(e_{i}\wedge e_{j})=(Te_{i})\wedge (Te_{j})
\end{equation}

The matrix $A$ transforms as: 
\begin{equation}
A\mapsto M^{T}AM
\end{equation}

We can write $\omega^k \in \Lambda^{2k} V$ as: 
\begin{equation}
\omega^k = \sum_I \alpha_I \, e_{i_1} \wedge \cdots \wedge e_{i_{2k}}, \quad 
\text{where } \alpha_I = \func{Pf}(A_I)
\end{equation}

Under the basis change $M$, the coefficients transform as: 
\begin{equation}
\vec{\alpha}^{\prime }= \mathrm{Pf}_k(M)^T \cdot \vec{\alpha}
\end{equation}

This parallels the role of compound matrices in determinant expansions.

\subsection*{2. Example: $4 \times 4$ Case for $\mathrm{Pf}_1(A)$}

Let: 
\begin{equation}
A = 
\begin{pmatrix}
0 & a & b & c \\ 
-a & 0 & d & e \\ 
-b & -d & 0 & f \\ 
-c & -e & -f & 0%
\end{pmatrix}%
\end{equation}

The 6 basis elements of $\Lambda^2 V$ are: 
\begin{align*}
e_{12}, e_{13}, e_{14}, e_{23}, e_{24}, e_{34}
\end{align*}

Label the Pfaffian compound matrix $\mathrm{Pf}_1(A)$ with rows/columns
indexed by these pairs. The diagonal entries are: 
\begin{align*}
\func{Pf}(A_{\{1,2\}}) &= a \\
\func{Pf}(A_{\{1,3\}}) &= b \\
\func{Pf}(A_{\{1,4\}}) &= c \\
\func{Pf}(A_{\{2,3\}}) &= d \\
\func{Pf}(A_{\{2,4\}}) &= e \\
\func{Pf}(A_{\{3,4\}}) &= f
\end{align*}

An off-diagonal entry: 
\begin{align*}
[\mathrm{Pf}_1(A)]_{\{1,2\},\{3,4\}} &= \func{Pf}(A_{\{1,2,3,4\}}) = af - be
+ cd
\end{align*}

Hence, the partial structure of $\mathrm{Pf}_1(A)$ is: 
\begin{equation}
\begin{pmatrix}
a & * & * & * & * & af - be + cd \\* 
& b & * & * & * & * \\* 
& * & c & * & * & * \\* 
& * & * & d & * & * \\* 
& * & * & * & e & * \\ 
af - be + cd & * & * & * & * & f%
\end{pmatrix}%
\end{equation}

The remaining off-diagonal entries are Pfaffians of $4 \times 4$ submatrices
formed by union of index pairs.

\section*{Summary}

\begin{itemize}
\item Under a basis change $M$, the coefficients of $\omega ^{k}$ in $%
\Lambda ^{2k}V$ transform via the Pfaffian compound matrix $\mathrm{Pf}%
_{k}(M)$.

\item For $V=\mathbb{R}^{4}$, $\mathrm{Pf}_{1}(A)\in \mathbb{R}^{6\times 6}$
has diagonal entries given by $2\times 2$ Pfaffians, and off-diagonal
entries built from $4\times 4$ Pfaffians.

\item This structure governs the transformation of even-grade wedge powers
and relates to Pl\"{u}cker coordinates.
\end{itemize}

\bigskip

\section{Transformation of Geminal Products}

\begin{definition}
Let $T\in \mathcal{B}\left( \mathcal{H}^{2}\right) $ i.e. 
\begin{equation}
h=Tg,\quad h,g\in \mathcal{H}^{2}
\end{equation}%
then the induced transformation $T^{p}$   
\begin{equation}
T^{p}:\mathcal{H}^{2p}\rightarrow \mathcal{H}^{2p}
\end{equation}%
is defined by 
\begin{equation}
T^{p}\left( g_{1}\wedge \cdots \wedge g_{p}\right) =T\left( g_{1}\right)
\wedge \cdots \wedge T\left( g_{p}\right) .
\end{equation}
\end{definition}

\begin{definition}
The associative product $\Theta \in \mathcal{B}\left( \mathcal{F}_{e}\right) 
$ 
\begin{equation}
T\overset{p}{\overbrace{\Theta \cdots \Theta }}T=T^{p},0\leq p\leq r.
\end{equation}
\end{definition}

\begin{theorem}
\begin{equation}
\mathcal{R}_{e}\left( T^{p}\right) =\func{Pf}_{p}\left( T\right) 
\end{equation}
\end{theorem}

\begin{proof}
\begin{eqnarray}
&&T^{p}\left( 
\begin{array}{ccc}
e_{12}\wedge \cdots \wedge e_{p-1p} & \cdots  & e_{\left( 2r-p\boldsymbol{-}%
1\right) \left( 2r-p\right) }\wedge \cdots \wedge e_{\left( 2r-1\right)
\left( 2r\right) }%
\end{array}%
\right)  \\
&=&\left( 
\begin{array}{ccc}
T\left( e_{12}\right) \wedge \cdots \wedge T\left( e_{p-1p}\right)  & \cdots 
& T\left( e_{\left( 2r-p\boldsymbol{-}1\right) \left( 2r-p\right) }\right)
\wedge \cdots \wedge T\left( e_{\left( 2r-1\right) \left( 2r\right) }\right) 
\end{array}%
\right)  \\
&=&\left( 
\begin{array}{ccc}
e_{12}\wedge \cdots \wedge e_{p-1p} & \cdots  & e_{\left( 2r-p\boldsymbol{-}%
1\right) \left( 2r-p\right) }\wedge \cdots \wedge e_{\left( 2r-1\right)
\left( 2r\right) }%
\end{array}%
\right) \mathcal{R}_{e}\left( T^{p}\right)  \\
&=&\left( 
\begin{array}{ccc}
e_{12}\wedge \cdots \wedge e_{p-1p} & \cdots  & e_{\left( 2r-p\boldsymbol{-}%
1\right) \left( 2r-p\right) }\wedge \cdots \wedge e_{\left( 2r-1\right)
\left( 2r\right) }%
\end{array}%
\right) \func{Pf}_{p}\left( T\right) 
\end{eqnarray}
\end{proof}

\begin{example}
\begin{eqnarray}
p &=&2,r=4 \\
2r &=&8 \\
2r-1 &=&7 \\
\left( 2r-\left( p\boldsymbol{+}1\right) \right)  &=&5 \\
\left( 2r-p\right)  &=&6
\end{eqnarray}%
\begin{equation}
\left( 
\begin{array}{ccc}
e_{12}\wedge \cdots \wedge e_{p-1p} & \cdots  & e_{\left( 2r-p\boldsymbol{-}%
1\right) \left( 2r-p\right) }\wedge \cdots \wedge e_{\left( 2r-1\right)
\left( 2r\right) }%
\end{array}%
\right) =\left( 
\begin{array}{ccc}
e_{12}\wedge e_{34} & \cdots  & e_{56}\wedge e_{78}%
\end{array}%
\right) 
\end{equation}
\end{example}

\end{document}
